\chapter{Guía rápida de montaje y uso}
\label{ap:guia-montaje-uso}

Este apéndice resume de forma esquemática las indicaciones prácticas para preparar el montaje físico y utilizar la aplicación en una partida asistida por visión.

\begin{figure}[H]
    \centering
    % \includegraphics[width=\textwidth]{imagenes/montaje_uso.png}
    \caption{Montaje físico recomendado para el uso del sistema.}
    \label{fig:montaje-fisico}
\end{figure}

\section{Montaje inicial del sistema}

\begin{enumerate}
    \item Colocar el tapete de Onitama dentro del encuadre de la cámara.

    \item Orientar el tapete de forma que, desde la perspectiva de la imagen, el lado azul quede a la izquierda y el lado rojo a la derecha.

    \item Reservar espacio dentro del encuadre para las cinco cartas visibles de la partida.

    \item En el montaje utilizado, la webcam del portátil se situó aproximadamente a 9 cm del borde más cercano del tapete y a unos 12 cm de altura.

    \item Ajustar ligeramente la altura de la cámara si es necesario para ampliar el campo de visión, siempre que las piezas sigan observándose con suficiente detalle.
\end{enumerate}

\section{Preparación de una partida física}

\begin{enumerate}
    \item Colocar las piezas en la configuración inicial de Onitama.

    \item Barajar las cartas de movimiento y seleccionar las cinco cartas visibles de la partida.

    \item Colocar las dos cartas del jugador rojo y las dos cartas del jugador azul en sus regiones correspondientes.

    \item Situar la carta lateral en una región fija, ubicada en el lado opuesto a la cámara.

    \item Colocar la carta lateral perpendicular al lado del tablero y orientarla hacia el jugador correspondiente.
\end{enumerate}

\section{Configuración y calibración}

\begin{enumerate}
    \item Abrir la aplicación gráfica del sistema desde un terminal mediante el comando \texttt{onitama-vision}.

    \item Seleccionar el color del jugador humano y el nivel de dificultad de la inteligencia artificial.

    \item Si la aplicación dispone de una calibración previa del tablero y de las regiones de cartas, y el montaje físico no ha cambiado, reutilizar dicha configuración.

    \item Si no existe una calibración previa, o si se ha modificado la posición de la cámara o el tapete, realizar de nuevo la calibración procurando mantener una disposición similar al montaje de referencia empleado durante el desarrollo:
    \begin{enumerate}[label=\alph*)]
        \item Para calibrar el tablero, marcar sus cuatro esquinas sobre la imagen capturada.

        \item Para calibrar las cartas, seleccionar cada carta mediante su botón correspondiente y marcar los cuatro puntos que delimitan su región de interés.
    \end{enumerate}

    \item Iniciar la partida desde la interfaz gráfica.
\end{enumerate}

\section{Recomendaciones durante el uso}

\begin{itemize}
    \item Mantener condiciones de iluminación estables.

    \item Evitar reflejos sobre las cartas y sombras intensas sobre el tablero o las piezas.

    \item Evitar oclusiones prolongadas de piezas o cartas.

    \item No mover la cámara ni el tapete durante la partida.

    \item Si la aplicación detecta una discrepancia, corregir la escena física antes de continuar.
\end{itemize}